\cotefile{32}
\scanpage{7}
% archive sheet number: 6

4. \uline{Homomorphisme de Gysin}.

Considérons un diagramme commutatif
\noindent(4.1)
\diagram{\begin{tikzcd}[ampersand replacement=\&,column sep=1.1em,row sep=1.5em]
X \handwritten{\arrow[rr,"f"]} \handwritten{\arrow[dr,"p"']} \& \& Y \handwritten{\arrow[dl,"q"]} \\
\& S \&
\end{tikzcd}},
avec $p$ et $q$ compactifiables, \inserted{$p$ plat de prés. finie, \deleted{\unsure{de dim}} $q$ lisse,} $\unsure{\deleted{A}S|\deleted{S}A}$ étant muni d'un Anneau de torsion $A$ premier aux car. résiduelles. Pour un complexe de $A$-modules $F$ sur $S$, on définit alors un homomorphisme canonique
\[
\handwritten{\mathrm{Tr}_f\ \text{ou}\ }f_{\handwritten{*}}:\quad \mathrm{R}f_{\handwritten{!}}(F_X)\handwritten{\longrightarrow} F_Y(d'-d)\handwritten{[}2(d'-d)\handwritten{]},
\tag*{(4.2)}
\]
en supposant $p,q$ de dimensions relatives $d,d'$ en tout point (il suffit que $p$ soit de dim rel $\handwritten{\leqslant} d$ en tout point). Si $p$ est également lisse et de dim. rel. exactement $d$ partout, alors l\unsure{\deleted{h}'|\deleted{'}h}homomorphisme précédent s'identifie à l'homomorphisme d'adjonction pour $\mathrm{R}f_{\handwritten{!}}$, $\mathrm{R}f^{\handwritten{!}}$ relatif au deuxième membre. En tous cas, l'homomorphisme en question s'appelle \uline{homomorphisme de Gysin} (relatif à $f,F$), ou \uline{homomorphisme trace} (\unsure{parfois} noté $\mathrm{Tr}_f$). On \unsure{\deleted{d}a|\deleted{a}d}\unsure{\deleted{a}p|\deleted{p}a}pelle encore par le même nom, et on note de même, tout homomorphisme déduit de (4.2) par application d'un foncteur quelconque (le \deleted{cas échéant}\inserted{\handwritten{plus souvent,}} en transformant encore la flèche obtenue en appliquant des isomorphismes canoniques, tels des isomorphismes de transitivité ... , utilisés communément \deleted{\unsure{à d}} comme morphismes d'identification ...), Par exemple on a les homomorphismes
\[
\mathrm{H}^i(f_*)\ \text{ou}\ f_*:\quad \mathbb{R}^i f_{\handwritten{!}}(F_X)\handwritten{\longrightarrow} \underline{\unsure{\mathbb{H}|\mathrm{H}}}^{i\inserted{\handwritten{+2(d'-d)}}}_{\deleted{\unsure{Y|*}}}(F_{\inserted{\handwritten{Y}}})_{\deleted{\unsure{Y|*}}}(d'-d)\deleted{\unsure{[2(d'-d)}},
\tag*{(4.2.1)}
\]
de faisceaux sur $Y$, déduit\handwritten{s} de (4.2?1) en appliquant $\underline{\mathrm{H}}^i$, ou
\[
\mathrm{R}q_{\handwritten{!}}(f_{\handwritten{*}})\ \text{ou}\ f_{\handwritten{*}}:\quad \mathrm{R}p_{\handwritten{!}}(F_X)\handwritten{\longrightarrow}\mathrm{R}q_{\handwritten{!}}(F_Y)(d'-d)\handwritten{[}2(d'-d)\handwritten{]},
\tag*{(4.2.2)}
\]
déduit de (4.2) en appliquant $\mathrm{R}q_{\handwritten{!}}$, ou plus généralement, pour tout morphisme com\unsure{\deleted{â}p|\deleted{p}â}actifiable
\[
g:\ Y\handwritten{\longrightarrow}Z
\]
